% AA-50 — Trzeci przypadek Kovacica i skończone grupy Galois.
% Wyniki klasyczne oznaczone jawnie: lokalne listy, kompletność
% i klasyfikacja prymitywnych grup tetra/octa/icosahedralnych.
% Udowodnione wprost: rekurencja i tożsamość przepływu Riccatiego;
% odzyskanie pierwiastka, końcowy test i przykład czworościenny.
\chapter{Trzeci przypadek algorytmu Kovacica i skończone grupy Galois}
\chaptermark{Trzeci przypadek Kovacica}
\label{aa:ch:kovacic-finite-groups}

W rozdziałach~\ref{aa:ch:kovacic-rational-potential}
i~\ref{aa:ch:kovacic-quadratic} poznaliśmy dwa
pierwsze przypadki algorytmu Kovacica.
Pierwszy wykrywa rozwiązania wymierne
równania Riccatiego $w'+w^2=q$,
drugi --- rozwiązania algebraiczne stopnia $2$
nad $\mathbb C(x)$. Trzeci przypadek,
związany ze skończonymi prymitywnymi
grupami Galois, wymaga osobnej konstrukcji.
Jej zasadniczym narzędziem jest
rekurencja prowadząca do wielomianu
$F(x,T)$, którego pierwiastki są
rozwiązaniami Riccatiego.

Rozdzielamy \emph{dowody bezpośrednie}
rekurencji, jej równoważności z
równaniem transportowym i
odtwarzania pierwiastków od
\emph{cytowanych twierdzeń klasycznych}
o wyczerpujących listach danych lokalnych
oraz skończonych grupach Galois.
Przykłady i testy dokładne weryfikują
konkretne konstrukcje, a nie zastępują
ogólnej klasyfikacji.

\section{Postać normalna i warunki wstępne}

Niech $K=\mathbb C(x)$ z derywacją
$D=\mathrm d/\mathrm dx$.
Rozpatrujemy równanie $z''=qz$
dla $q\in K$. Zakładamy
efektywną reprezentację współczynników
w sensie
Uwagi~\ref{aa:uw:kov-rational-effective};
w szczególności dysponujemy dokładnymi
działaniami i rozkładem wielomianów
na czynniki w odpowiednim
algebraicznym rozszerzeniu.

\begin{definicja}[Rząd potencjału w nieskończoności]
\label{aa:def:kov3-infinity-order}
Niech $q=a/b\in\mathbb C(x)\setminus\{0\}$,
gdzie $a,b\in\mathbb C[x]$ są względnie
pierwsze i $b\ne0$.
Kładziemy
$\operatorname{ord}_{\infty}q=\deg b-\deg a$.
Jeśli $\operatorname{ord}_{\infty}q=2$,
przez $b_\infty$ oznaczamy współczynnik
przy $x^{-2}$ w rozwinięciu
Laurenta $q$ przy nieskończoności.
Przy skończonym biegunie $c$ rzędu $2$
oznaczamy
$b_c=[(x-c)^{-2}]_c q$.
\end{definicja}

\begin{definicja}[Prymitywna reprezentacja wymiaru drugiego]
\label{aa:def:kov3-primitive}
Niech $G$ będzie grupą,
a $\rho:G\to\operatorname{GL}(V)$
reprezentacją zespoloną na
dwuwymiarowej przestrzeni $V$.
Mówimy, że jest ona
\emph{prymitywna}, jeśli
jest nieprzywiedlna oraz
nie istnieją różne
jednowymiarowe podprzestrzenie
$L_1,L_2\subset V$ takie, że
$V=L_1\oplus L_2$ i
zbiór $\{L_1,L_2\}$
jest zachowywany przez
$\rho(G)$.
Jeżeli taki rozkład
istnieje, reprezentację
nieprzywiedlną nazywamy
\emph{imprymitywną}.
\end{definicja}

\begin{twierdzenie}[Przypadek trzeci: klasyfikacja -- wynik klasyczny]
\label{aa:tw:kov3-classification}
Niech $q\in\mathbb C(x)$.
Jeśli różniczkowa grupa Galois
równania $z''=qz$ nad $\mathbb C(x)$
jest skończona i jej reprezentacja
wymiaru $2$ jest prymitywna,
to jej obraz w
$\operatorname{PSL}_2(\mathbb C)$
jest izomorficzny z jedną z grup
$A_4$, $S_4$ lub $A_5$.
Odpowiadające im parametry
trzeciego przypadku algorytmu Kovacica
należą do zbioru
\[
 n\in\{4,6,12\}.
\]
W tym przypadku wszystkie
rozwiązania równania liniowego
są algebraiczne nad $\mathbb C(x)$.
Wyczerpujący charakter tej klasyfikacji
jest klasycznym wynikiem teorii
Picarda--Vessiota i algorytmu Kovacica.
\end{twierdzenie}

\begin{uwaga}[Źródło]
\label{aa:uw:kov3-classical}
Korzystamy z pracy J.~Kovacica,
\emph{An algorithm for solving second
order linear homogeneous differential
equations}, \emph{Journal of Symbolic
Computation} 2 (1986), 3--43.
W szczególności nie twierdzimy, że
poniższe rachunki z jednym przykładem
dowodzą klasyfikacji skończonych
podgrup $\operatorname{PSL}_2(\mathbb C)$
lub pełnej odpowiedniości Galois.
Pierwszy i drugi przypadek algorytmu
zostały omówione wcześniej;
sam wynik negatywny trzeciego przypadku
jest konkluzją dla całego algorytmu
dopiero po wykluczeniu obu pierwszych.
\end{uwaga}

\section{Dane lokalne trzeciego przypadku}

\begin{twierdzenie}[Wyczerpujące listy lokalne -- wynik klasyczny]
\label{aa:tw:kov3-local-sets}
Niech $q\in K\setminus\{0\}$.
W trzecim przypadku Kovacica każdy
skończony biegun $q$ ma rząd $1$ albo $2$
oraz $\operatorname{ord}_{\infty}q\ge2$.
Oznaczmy przez $\Gamma$ zbiór
wszystkich skończonych biegunów $q$.
Dla każdego $n\in\{4,6,12\}$
definiujemy następujące
skończone zbiory liczb całkowitych:
\begin{equation}\label{aa:eq:kov3-local-sets}
 E_c(n)=
 \begin{cases}
 \{12\},&\text{jeśli biegun w $c$ jest prosty},\\[2pt]
 \left\{6+\dfrac{12k}{n}\sqrt{1+4b_c}:
       k=-\dfrac n2,\ldots,\dfrac n2\right\}
       \cap\mathbb Z,
       &\text{jeśli ma rząd $2$};
 \end{cases}
\end{equation}
przy nieskończoności przyjmujemy
\begin{equation}\label{aa:eq:kov3-infinity-sets}
 E_\infty(n)=
 \begin{cases}
 \left\{6+\dfrac{12k}{n}:
       k=-\dfrac n2,\ldots,\dfrac n2\right\}
       \cap\mathbb Z,
       &\operatorname{ord}_{\infty}q>2,\\[2pt]
 \left\{6+\dfrac{12k}{n}\sqrt{1+4b_\infty}:
       k=-\dfrac n2,\ldots,\dfrac n2\right\}
       \cap\mathbb Z,
       &\operatorname{ord}_{\infty}q=2 .
 \end{cases}
\end{equation}
Jeśli występuje skończona
prymitywna grupa Galois, to
dla pewnego $n\in\{4,6,12\}$
i pewnego wyboru
$e_c\in E_c(n)$,
$e_\infty\in E_\infty(n)$,
liczba
\begin{equation}\label{aa:eq:kov3-degree}
 d=\frac n{12}
       \left(e_\infty-\sum_{c\in\Gamma}e_c\right)
\end{equation}
jest nieujemną liczbą całkowitą.
Dla takiego wyboru definiujemy
\begin{equation}\label{aa:eq:kov3-theta-S}
 \theta=\frac n{12}\sum_{c\in\Gamma}
                   \frac{e_c}{x-c},
 \qquad S=\prod_{c\in\Gamma}(x-c).
\end{equation}
W klasycznej procedurze
istnieje unormowany
$P\in\mathbb C[x]$ stopnia $d$
spełniający podaną poniżej
rekurencję z warunkiem
$P_{-1}=0$.
\end{twierdzenie}

\begin{uwaga}[Dokładna granica cytatu]
\label{aa:uw:kov3-sets-cited}
Twierdzenie~\ref{aa:tw:kov3-local-sets}
przyjmujemy jako klasyczną,
wyczerpującą część algorytmu
Kovacica. Samodzielnie
nie dowodzimy ani kompletności
list $E_c(n)$ i $E_\infty(n)$,
ani wynikającej z teorii
niezmienników konieczności
istnienia wielomianu $P$.
Natomiast \emph{poprawność
każdego znalezionego wyniku}
wyniknie z dowodu równania
transportowego. Jeśli $q=0$,
rozwiązania $z=1,x$ są już
znane i nie uruchamiamy
przypadku trzeciego.
\end{uwaga}

\section{Rekurencja wielomianowa}

\begin{definicja}[Wielomiany pomocnicze]
\label{aa:def:kov3-recurrence}
Ustalmy $n\in\{4,6,12\}$,
dane $S,\theta$ z
\eqref{aa:eq:kov3-theta-S}
oraz unormowany wielomian
$P$ stopnia $d$.
Definiujemy
$P_{n+1}=0$, $P_n=-P$
i rekurencyjnie dla
$i=n,n-1,\ldots,0$:
\begin{equation}\label{aa:eq:kov3-recurrence}
 P_{i-1}
 =-SP_i'
 +\bigl((n-i)S'-S\theta\bigr)P_i
 -(n-i)(i+1)S^2qP_{i+1}.
\end{equation}
Każdy $P_i$ jest wielomianem,
ponieważ $S\theta$ i
$S^2q$ są wielomianami
przy założeniach
Twierdzenia~\ref{aa:tw:kov3-local-sets}.
Warunkiem końcowym jest
$P_{-1}=0$ jako tożsamość
w $\mathbb C[x]$.
\end{definicja}

Wbrew pozorom indeks $i=-1$
nie oznacza pierwiastka
równania różniczkowego.
Jest to warunek zgodności
ostatniego kroku rekurencji,
a jego współczynniki
dają skończony układ
liniowy na nieznane
współczynniki $P$.

\begin{lemat}[Rekurencja jako równanie transportowe]
\label{aa:lem:kov3-transport}
Niech $q,S,\theta$ będą
funkcjami wymiernymi,
$S\ne0$, a $n\ge1$.
Niech funkcje $P_i\in K$
spełniają
\eqref{aa:eq:kov3-recurrence}
dla $i=n,\ldots,0$,
z $P_{n+1}=0$, $P_n=-P\ne0$
oraz $P_{-1}=0$.
Wprowadźmy
\begin{equation}\label{aa:eq:kov3-polynomial-F}
 F(x,T)=\sum_{i=0}^{n}
          \frac{S^iP_i}{(n-i)!}T^i\in K[T]
\end{equation}
i derywację
$\mathcal D=\partial_x+(q-T^2)\partial_T$
na $K[T]$.
Wtedy
\begin{equation}\label{aa:eq:kov3-transport}
 \mathcal DF
 =\left(-nT+n\frac{S'}S-\theta\right)F.
\end{equation}
Odwrotnie, przy
$P_{n+1}=0$, $P_n=-P\ne0$
tożsamość~\eqref{aa:eq:kov3-transport}
jest równoważna rekurencji
i warunkowi $P_{-1}=0$.
\end{lemat}
\begin{proof}
Oznaczmy
$A_i=S^iP_i/(n-i)!$ dla
$0\le i\le n$, a poza tym
zakresem $A_{-1}=A_{n+1}=0$.
Współczynnik przy $T^i$
w $\mathcal DF$ dla
$0\le i\le n$ wynosi
\[
 A_i'+(i+1)qA_{i+1}-(i-1)A_{i-1}.
\]
Współczynnik przy tej samej
potędze w prawej stronie
\eqref{aa:eq:kov3-transport}
wynosi
$-nA_{i-1}+
 (nS'/S-\theta)A_i$.
Ich równość można
przekształcić do
\[
 (n-i+1)A_{i-1}
 =-A_i'+\left(n\frac{S'}S-\theta\right)A_i
 -(i+1)qA_{i+1}.
\]
Po podstawieniu $A_i$
i pomnożeniu przez
$(n-i)!/S^{i-1}$
otrzymujemy dokładnie
\eqref{aa:eq:kov3-recurrence}.
Dla $i=0$ powstaje
warunek $P_{-1}=0$.
Przy $T^{n+1}$ po obu
stronach otrzymujemy
$-nA_n$, więc nie ma
dodatkowego warunku.
To dowodzi także
równoważności w drugą stronę.
\end{proof}

\begin{lemat}[Algebraiczne pierwiastki rozwiązują Riccatiego]
\label{aa:lem:kov3-roots}
Niech $F\in K[T]\setminus K$
oraz $H\in K[T]$
spełniają $\mathcal DF=HF$,
gdzie
$\mathcal D=\partial_x+(q-T^2)\partial_T$
i $q\in K$.
Wtedy każdy pierwiastek
$w$ dowolnego nierozkładalnego
czynnika wielomianu $F$
nad $K$ spełnia
$w'+w^2=q$ w jedynym
przedłużeniu derywacji
do $K(w)$.
W szczególności otrzymujemy
algebraiczne rozwiązanie
równania Riccatiego.
\end{lemat}
\begin{proof}
Rozłóżmy w $K[T]$
wielomian $F$ na
niezerowe nierozkładalne
czynniki $f_j$ o
krotnościach $m_j$.
Z tożsamości
$\mathcal DF=HF$ wynika
\[
 \sum_jm_j\frac{\mathcal Df_j}{f_j}=H.
\]
W pierścieniu ilorazowym
$K[T]_{(f_j)}$ pozostałe
czynniki są odwracalne.
Ponieważ prawa strona
jest wielomianem,
wyraz $\mathcal Df_j/f_j$
nie może mieć bieguna,
czyli $f_j$ dzieli
$\mathcal Df_j$.
Niech $w$ będzie
pierwiastkiem $f_j$.
W charakterystyce $0$
nierozkładalny $f_j$
jest separowalny, a więc
$\partial_T f_j(x,w)\ne0$.
Zróżniczkowanie
$f_j(x,w)=0$ daje
\[
 0=\partial_x f_j(x,w)
   +\partial_T f_j(x,w)w'.
\]
Jednocześnie
$\mathcal D f_j(x,w)=0$,
a zatem
$\partial_x f_j(x,w)
 +(q-w^2)\partial_T f_j(x,w)=0$.
Odjęcie tych równości
i podzielenie przez
$\partial_Tf_j(x,w)$
daje $w'=q-w^2$.
\end{proof}

\begin{twierdzenie}[Skończona procedura i jej poprawność]
\label{aa:tw:kov3-algorithm}
Niech $q\in\mathbb C(x)$
ma efektywnie zadane
współczynniki.
Przyjmując klasyczne
Twierdzenie~\ref{aa:tw:kov3-local-sets}
i wcześniejszą klasyfikację
Kovacica, trzecia gałąź
jest rozstrzygalna:
\begin{enumerate}
\item odrzucamy $q\ne0$
o skończonym biegunie
rzędu większego niż $2$
lub o
$\operatorname{ord}_{\infty}q<2$;
\item dla kolejnych
$n=4,6,12$ obliczamy
skończone zbiory
$E_c(n)$ i $E_\infty(n)$;
\item rozpatrujemy każdą
kombinację, dla której
$d$ z~\eqref{aa:eq:kov3-degree}
jest nieujemną liczbą
całkowitą, i tworzymy
$S,\theta$;
\item wstawiamy
$P=x^d+\sum_{j=0}^{d-1}p_jx^j$
do~\eqref{aa:eq:kov3-recurrence}
i rozwiązujemy
skończony liniowy układ
powstały z $P_{-1}=0$;
\item dla znalezionego
$P$ tworzymy $F$
z~\eqref{aa:eq:kov3-polynomial-F},
znajdujemy jego
pierwiastek algebraiczny
$w$ i otrzymujemy
lokalne niezerowe rozwiązanie
$z=\exp(\int w\,\mathrm dx)$.
\end{enumerate}
Odpowiedź dodatnia jest
weryfikowalna bez odwołania
do klasyfikacji grup Galois.
Jeśli pierwsze dwie gałęzie
Kovacica zostały wykluczone,
negatywny wynik wszystkich
trzech gałęzi oznacza brak
niezerowych rozwiązań
liouvillowskich.
\end{twierdzenie}
\begin{proof}
Zbiorów danych lokalnych
jest skończenie wiele.
Dla każdej kombinacji
znany jest stopień $d$,
a rekurencja jest liniowa
w $P$ i jego pochodnych.
Ponieważ $S\theta$ i
$S^2q$ są wielomianami,
każdy $P_i$ jest
wielomianem liniowo
zależnym od nieznanych
$p_j$, więc test
$P_{-1}=0$ sprowadza się
do skończonego układu
równań liniowych.
Dla dowolnego dodatniego
wyniku konstrukcja $F$
spełnia równanie
\eqref{aa:eq:kov3-transport}
na mocy
Lematu~\ref{aa:lem:kov3-transport}.
Ponieważ współczynnik
wiodący $F$ przy $T^n$
wynosi $-S^nP\ne0$,
wielomian nie jest stały
w $T$. Lemma
\ref{aa:lem:kov3-roots}
daje pierwiastek
$w$ spełniający Riccatiego.
Stąd $z' = wz$ prowadzi do
$z''=qz$.
Znalezienie pierwiastka
w skończonym rozszerzeniu
algebraicznym jest
dopuszczalne w zadanym
efektywnym modelu.
Natomiast \emph{konieczność}
istnienia jednej z
kombinacji i wielomianu $P$
dla skończonej prymitywnej
grupy Galois wynika
właśnie z cytowanego
Twierdzenia~\ref{aa:tw:kov3-local-sets}.
Ostatni wniosek o braku
wszystkich rozwiązań
liouvillowskich wykorzystuje
pełną klasyfikację czterech
przypadków Kovacica
z~\ref{aa:tw:kovacic-classification},
również cytowaną
jako wynik klasyczny.
\end{proof}

\section{Przykład czworościenny}

\begin{przyklad}[Potencjał typu czworościennego]
\label{aa:prz:kov3-tetrahedral}
Rozważmy
\begin{equation}\label{aa:eq:kov3-tetra-q}
 q=-\frac3{16x^2}
    -\frac2{9(x-1)^2}
    +\frac3{16x(x-1)}.
\end{equation}
Skończone bieguny
w $0$ i $1$ mają rząd $2$,
a w nieskończoności
$\operatorname{ord}_\infty q=2$.
Współczynniki wiodące są
\[
 b_0=-\frac3{16},\qquad
 b_1=-\frac29,\qquad
 b_\infty=-\frac29.
\]
Przy $n=4$ pierwiastki
$\sqrt{1+4b_c}$ wynoszą
odpowiednio $1/2$, $1/3$, $1/3$.
Zatem
\[
 E_0(4)=\{3,6,9\},\qquad
 E_1(4)=E_\infty(4)=\{4,5,6,7,8\}.
\]
Wybór
$e_0=3$, $e_1=4$,
$e_\infty=7$ daje
\[
 d=\frac13(7-3-4)=0,\qquad
 P=1,\quad
 S=x(x-1),\quad
 \theta=\frac1x+\frac4{3(x-1)}.
\]
Rekurencja~\eqref{aa:eq:kov3-recurrence}
przy $n=4$ wyznacza
\[
\begin{aligned}
 P_4&=-1,\\
 P_3&=\frac{7x-3}{3},\\
 P_2&=-\frac{48x^2-41x+9}{12},\\
 P_1&=\frac{320x^3-409x^2+180x-27}{72},\\
 P_0&=-\frac{2048x^4-3484x^3+2313x^2
                -702x+81}{864},\\
 P_{-1}&=0.
\end{aligned}
\]
Tak więc wielomian
\[
 F(x,T)=\sum_{i=0}^{4}
       \frac{S^iP_i}{(4-i)!}T^i
\]
ma stopień $4$ względem $T$,
a każdy pierwiastek
nierozkładalnego czynnika
$F$ nad $\mathbb C(x)$
spełnia $w'+w^2=q$.
Otrzymaliśmy w pełni
weryfikowalny dodatni wynik
procedury bez zapisywania
rozwiniętego wielomianu
o dużych współczynnikach.
Przykład odpowiada
lokalnym różnicom wykładników
$1/2$, $1/3$, $1/3$,
typowym dla grupy czworościanu.
Przypisanie pełnej grupy
Galois korzysta z klasycznej
teorii monodromii
równania hipergeometrycznego.
\end{przyklad}

\begin{przyklad}[Przypadki ośmiościenny i dwudziestościenny]
\label{aa:prz:kov3-octa-icosa}
Dwa kolejne przykłady
mają wspólną postać
\[
 q_m=-\frac3{16x^2}
     -\frac2{9(x-1)^2}
     +\frac{C_m}{x(x-1)},\qquad
 C_4=\frac{101}{576},\quad
 C_5=\frac{611}{3600}.
\]
Dla $m=4$ odpowiadające
różnice wykładników lokalnych
przy $0,1,\infty$ wynoszą
$1/2,1/3,1/4$, czyli
dane typu ośmiościennego.
Przy $n=6$ zbiorami
kandydatów są
\[
 E_0=\{3,4,5,6,7,8,9\},
 \quad E_1=\{4,6,8\},
 \quad E_\infty=\{5,6,7\}.
\]
Wybór $e_0=3$, $e_1=4$,
$e_\infty=7$ daje $d=0$,
$P=1$ i
\[
 S=x(x-1),\qquad
 \theta=\frac3{2x}+\frac2{x-1},
 \qquad S\theta=\frac72x-\frac32,
\]
a także
\[
 S^2q_4
 =-\frac{15}{64}x^2
  +\frac{115}{576}x-\frac3{16}.
\]
Wykonanie skończonej
rekurencji
\eqref{aa:eq:kov3-recurrence}
dla $i=6,\ldots,0$
prowadzi dokładnie
do $P_{-1}=0$.
Współczynniki są wymierne,
zatem wszystkie etapy można
sprawdzić arytmetycznie,
bez przybliżeń.

Dla $m=5$ otrzymujemy
różnice wykładników
$1/2,1/3,1/5$ i dane
typu dwudziestościennego.
Przy $n=12$ mamy
\[
 E_0=\{3,4,5,6,7,8,9\},
 \quad E_1=\{4,5,6,7,8\},
 \quad E_\infty=\{5,6,7\}.
\]
Ponownie $e_0=3$, $e_1=4$,
$e_\infty=7$ daje $d=0$,
$P=1$, lecz tym razem
\[
 \theta=\frac3x+\frac4{x-1},
 \qquad S\theta=7x-3,
\]
oraz
\[
 S^2q_5
 =-\frac6{25}x^2
  +\frac{739}{3600}x-\frac3{16}.
\]
Rekurencja dla
$i=12,\ldots,0$ również
daje $P_{-1}=0$.
Równości dla obu wartości
$n$ są dodatkowo sprawdzane
przez niezależne testy
symboliczne i wymierne.
Przyporządkowanie
pełnej grupy monodromii
typom ośmio- i
dwudziestościennemu
wykorzystuje klasyczną
klasyfikację Schwarzowską,
a nie sam rachunek
$P_{-1}=0$.
\end{przyklad}

\begin{przyklad}[Rola dokładnej liczby całkowitej]
\label{aa:prz:kov3-integer-sets}
Dla $q$ z~\eqref{aa:eq:kov3-tetra-q}
i $n=4$ współczynnik
$b_0=-3/16$ daje
$\sqrt{1+4b_0}=1/2$.
W zbiorze
$\{6+3k/2:k=-2,-1,0,1,2\}$
nie wszystkie liczby są
całkowite:
$\{3,9/2,6,15/2,9\}$.
Lista $E_0(4)$ wynosi
zatem dokładnie $\{3,6,9\}$.
Niedopuszczalne jest
zaokrąglanie $9/2$ czy
$15/2$ do liczb całkowitych.
Podobnie w przypadku
$b_1=-2/9$ pierwiastek
wynosi $1/3$, a
$\{6+k:k=-2,\ldots,2\}$
daje pięć różnych
wartości całkowitych.
\end{przyklad}

\begin{przyklad}[Wcześniejsza gałąź nie znika]
\label{aa:prz:kov3-rational-degenerate}
Dla $q=2/x^2$ już
$w=2/x$ i $w=-1/x$
są rozwiązaniami
\emph{wymiernymi}
Riccatiego, ponieważ
$w'+w^2=q$.
Przypadek trzeci
nie jest tutaj potrzebny.
Ewentualne rozwiązanie
rekurencji~\eqref{aa:eq:kov3-recurrence}
nie upoważniałoby do
przypisania grupie Galois
typu czworościennego:
najpierw należy odróżnić
przypadki wymierny
i kwadratowy od skończonego
prymitywnego.
\end{przyklad}

\begin{przyklad}[Natychmiastowa przeszkoda]
\label{aa:prz:kov3-airy}
Dla równania Airy'ego
$z''=xz$ potencjał
$q=x$ nie ma skończonych
biegunów, ale
$\operatorname{ord}_{\infty}q=-1<2$.
Trzeci przypadek jest więc
wykluczony już przez
warunek wstępny z
Twierdzenia~\ref{aa:tw:kov3-local-sets}.
Łącznie z wynikami
rozdziałów~\ref{aa:ch:kovacic-rational-potential}
i~\ref{aa:ch:kovacic-quadratic}
oraz klasyczną klasyfikacją
Kovacica wyklucza to
liouvillowskość,
zgodnie z analizą
grupy Galois Airy'ego
w~\ref{aa:ch:picard-vessiot}.
\end{przyklad}

\section{Zamknięcie trzech przypadków}

Mamy teraz trzy osobne
procedury wraz z ich
granicami:
pierwszy przypadek bada
$w\in\mathbb C(x)$;
drugi --- $w$ algebraiczne
stopnia $2$;
trzeci --- przypadek
skończonej prymitywnej
grupy Galois, którego
właściwe stopnie algebraiczne
to $4$, $6$ i $12$.
Twierdzenie o wyczerpaniu
tych możliwości i
wniosek o braku rozwiązań
liouvillowskich w czwartym
przypadku są wynikiem
klasycznym, a nie
następstwem skończonej
liczby testów regresyjnych.

W odróżnieniu od
deklaracji teoretycznej
o pełnym algorytmie Rischa,
wszystkie trzy gałęzie
Kovacica są tu
\emph{opisane algorytmicznie},
ale ich matematyczna
kompletność poza
bezpośrednio udowodnionymi
redukcjami wykorzystuje
cytowaną klasyczną
klasyfikację lokalną
i teorię Picarda--Vessiota.
